\section{Task-Type Coverage and Limitations}\label{app:coverage}

This section honestly maps HARPO's demonstrated capability onto the task
taxonomy defined by the FPT'26 Track~A (LLM4HLS) call --- a useful external yardstick
of agent completeness, independent of any competition entry. Coverage claims below
were checked against the code (\texttt{harpo/}), not memory; numbers and proofs
trace to the canonical results table. Throughout, status is reported with three
labels: \textbf{Proven} (\checkmark) --- demonstrated end-to-end in the code with a
fixture or result; \textbf{Partial} (\(\sim\)) --- a path exists but is unproven; and
\textbf{Gap} (\(\times\)) --- a genuine gap.

\subsection{Bottom line}

HARPO proves the \emph{complete required agent workflow} on the two most central
task types --- \emph{optimize a correct-but-unoptimized baseline} and \emph{repair a
csim failure} --- with the diagnosis/budget architecture already reaching toward the
others. We claim performance and resource, not power. The probe is deliberately
capped (no broad design-space exploration). The remaining items are a scoping list
(Sec.~\ref{app:coverage:future}); several depend on the interface of whatever
external evaluation harness the agent is pointed at, and are best built against a
concrete harness rather than a guessed one.

\subsection{Task initial conditions}

The taxonomy spans the task types in Table~\ref{tab:cov-initial}. The table maps
each task type onto the HARPO mechanism that addresses it and its proven
status.

\begin{table*}[t]
\centering
\small
\caption{Task initial conditions: task type vs.\ HARPO mechanism and status.}
\label{tab:cov-initial}
\begin{tabular}{@{}p{0.30\textwidth}p{0.42\textwidth}p{0.22\textwidth}@{}}
\toprule
\textbf{Task type} & \textbf{HARPO mechanism} & \textbf{Status} \\
\midrule
Functionally correct but \textbf{unoptimized} baseline
& \texttt{run\_optimize} (5 hand-built + 3 PolyBench kernels)
& \checkmark~\textbf{Proven} \\
\addlinespace
\textbf{Fails compilation}
& \texttt{parser} $\rightarrow$ \texttt{compile\_error}; \texttt{diagnosis} $\rightarrow$ \texttt{minimal\_compile\_fix}; repair loop drives it
& \(\sim\)~Partial: path exists; \textbf{no fixture proves it} (repair fixtures are functional bugs, not compile failures) \\
\addlinespace
\textbf{Fails synthesis}
& \texttt{SYNTHESIS\_FAIL} diagnosis class + \texttt{parse\_csynth}, but \texttt{run\_repair} runs the \textbf{csim (gpp) stage only}
& \(\times\)~\textbf{Gap}: a design that passes csim yet fails csynth is not driven by any repair loop \\
\addlinespace
Compiles but \textbf{fails csim}
& \texttt{run\_repair} (\texttt{vadd\_buggy} / \texttt{offbyone} / \texttt{scale\_wrongop})
& \checkmark~\textbf{Proven} \\
\addlinespace
Compiles but \textbf{fails cosim}
& cosim is budgeted + stage-gated, but \texttt{runner} marks it ``(later)''
& \(\times\)~\textbf{Gap}: cosim not implemented \\
\addlinespace
Compiles but fails \textbf{hidden functional tests}
& agent runs the \emph{provided} testbench; hidden tests are evaluator-side
& \checkmark~Covered (by design --- nothing to build) \\
\addlinespace
\textbf{Structural: severe resource inefficiency}
& \texttt{run\_optimize} + \texttt{RESOURCE\_OVERUSE}
& \checkmark~Proven \\
\addlinespace
\textbf{Structural: deadlock}
& \texttt{TIMEOUT\_OR\_DEADLOCK} class + csim timeout detection
& \(\sim\)~Partial: class exists; \textbf{no streaming/dataflow kernel demonstrated} \\
\addlinespace
\textbf{Structural: invalid streaming behavior}
& --- (no \texttt{hls::stream}/dataflow kernel in the task set)
& \(\times\)~\textbf{Gap} \\
\addlinespace
Other HLS compilation problems
& catch-all via the \texttt{compile\_error} path
& \(\sim\)~Partial \\
\bottomrule
\end{tabular}
\end{table*}

\subsection{Required agent workflow (the 6 steps)}

The Track~A call specifies a six-step agent workflow. Table~\ref{tab:cov-workflow}
shows that HARPO covers all six. This row block is HARPO's strongest alignment ---
it was built around exactly this loop.

\begin{table*}[t]
\centering
\small
\caption{Required agent workflow: all six workflow steps are covered as mechanisms (task-type coverage gaps are in Tables~\ref{tab:cov-initial} and~\ref{tab:cov-artifacts}).}
\label{tab:cov-workflow}
\begin{tabular}{@{}p{0.34\textwidth}p{0.50\textwidth}p{0.08\textwidth}@{}}
\toprule
\textbf{Required step} & \textbf{HARPO} & \textbf{Status} \\
\midrule
1. Interpret the task spec + initial code
& \texttt{task.py} $\rightarrow$ \texttt{TaskContext} (part/clock/sources/policy/budget injected)
& \checkmark \\
\addlinespace
2. Generate/modify HLS C/C++ incl.\ pragmas
& \texttt{recipes.py} (deterministic) + \texttt{patch\_engine.py} (LLM)
& \checkmark \\
\addlinespace
3. Invoke tool feedback interfaces
& \texttt{runner.py} (gpp csim, vitis\_hls csynth)
& \checkmark \\
\addlinespace
4. Parse logs/reports for diagnosis
& \texttt{parser.py} + \texttt{diagnosis.py}
& \checkmark \\
\addlinespace
5. Correctness before QoR
& structural invariant: every candidate re-runs csim \emph{before} its csynth metrics are read; correctness tier dominates the score
& \checkmark \\
\addlinespace
6. Terminate within budget
& \texttt{budget.py} (\texttt{policy\_allows} checks \texttt{can()} before every call; a reserve is held back so exploration cannot spend the last verification call)
& \checkmark \\
\bottomrule
\end{tabular}
\end{table*}

\subsection{Provided task artifacts}

Table~\ref{tab:cov-artifacts} maps the artifacts the agent must consume against
HARPO's handling of each.

\begin{table*}[t]
\centering
\small
\caption{Provided task artifacts the agent must consume.}
\label{tab:cov-artifacts}
\begin{tabular}{@{}p{0.32\textwidth}p{0.40\textwidth}p{0.22\textwidth}@{}}
\toprule
\textbf{Artifact} & \textbf{HARPO} & \textbf{Status} \\
\midrule
Source files (baseline C/C++)
& injected via \texttt{TaskContext}
& \checkmark~Proven \\
\addlinespace
Testbench + \textbf{build scripts (e.g.\ Makefile)}
& HARPO uses its \textbf{own} runner/tcl, not a provided Makefile
& \(\sim\)~Partial: \textbf{adapter likely needed} to the evaluator's build/run interface \\
\addlinespace
Spec: interface contract, data types
& \texttt{spec.json} carries interface/types
& \checkmark~Proven \\
\addlinespace
Spec: \textbf{numerical tolerance}
& csim compare is exact today
& \(\sim\)~Partial: tolerance-aware compare may be needed \\
\addlinespace
Spec: design constraints
& spec/constraints JSON
& \checkmark~Proven \\
\addlinespace
Target constraints (part/clock/optional resource limits)
& injected; \texttt{RESOURCE\_OVERUSE} handles limits
& \checkmark~Proven \\
\addlinespace
\textbf{Budget config} --- max csim/cosim/synth \textbf{OR a unified credit budget}
& \texttt{BudgetManager} is \textbf{per-tool only} (the \texttt{mode} key exists but \texttt{can()} never branches on it)
& \(\times\)~\textbf{Gap}: unified-credit mode not implemented \\
\bottomrule
\end{tabular}
\end{table*}

\subsection{Threat to validity: csynth estimates are not measurements}\label{app:coverage:implverify}

Every number in this paper is a Vitis HLS csynth estimate. The repository also
carries a post-route fidelity check (\texttt{optimize --impl-verify 3},
2026-07-16; records under \texttt{docs/case-study/}): 12 candidates across 4 tasks
taken through \texttt{export\_design -flow impl} and out-of-context Vivado place
\& route on the same \texttt{xc7z020clg400-1} at 10\,ns, three of them
(\texttt{atax\_001}, \texttt{bicg\_001}, \texttt{gemm\_001}) kernels reported here.
LUT estimates are pessimistic in all 12 cases, by 1.60$\times$--2.73$\times$
(2.03$\times$--2.46$\times$ among winners). Clock is worse than estimated in 11 of
12: designs printed here at 144.45\,MHz measure 99.97--109.66\,MHz post-route. Two
candidates pass timing on the estimate and miss it after routing ---
\texttt{gemm\_001} cand\_0001 at 10.168\,ns (the worst degradation, $+3.256$\,ns)
and \texttt{atax\_001} cand\_0000 at 10.003\,ns. For the latter the defensible
claim is the $\approx$3\,ns estimate error, not a failing baseline: 0.003\,ns sits
inside plausible run-to-run resolution, place \& route ran once per candidate at
the default seed, and that variance was never measured.

Ranking diverges on 1 of 4 tasks: on \texttt{atax\_001} the estimates rank
cand\_0002 first, measurement cand\_0003. The first-place flip is
robust (109 LUT, 6.0\%); the 2nd/3rd pair differ by 4 LUT (0.2\%) and are
unordered. Measurement promotes a candidate this paper does not report.

A 1.6--2.7$\times$ spread does not threaten the direction of the headline area
result --- a $\approx$42$\times$ gap sits far outside it. It threatens marginal
verdicts. Sec.~\ref{app:scoring:comparison} calls \texttt{atax\_001}
``recipe (marginal)'' on an estimated 3907 vs.\ 3994 LUT, a 2.2\% margin ---
smaller than the measured estimate error on those very designs. That verdict is
not safe on estimates. HARPO optimizes an estimate ranking, and an estimate
ranking is not the silicon ranking.

The check is narrower than it looks. Even in the measured records, latency
and interval keep \texttt{latency\_source} = \texttt{csynth}: post-route cycle
counts were never measured. Area and timing are checked here; \emph{measured}
throughput is not.

\subsection{Future work}\label{app:coverage:future}

HARPO calls Vitis directly through its own \texttt{runner}; integrating with an
externally provided evaluation harness (one that supplies its own build scripts and
csim/cosim/synth entry points) would require an adapter, which is best written
against a concrete harness interface. In rough priority order:

\begin{enumerate}
  \item \textbf{Eval-interface adapter} --- wrap \texttt{runner} to call an external
    harness's csim/cosim/synth interfaces and consume its provided Makefile/build
    scripts.
  \item \textbf{Unified-credit budget mode} --- branch \texttt{BudgetManager.can()} on
    \texttt{mode} so a single credit pool is honored, if that's what the evaluator meters.
  \item \textbf{Synthesis-failure repair} --- drive \texttt{run\_repair} (or a unified
    repair stage) from \texttt{SYNTHESIS\_FAIL}/\texttt{TIMING\_FAIL}, not csim only.
  \item \textbf{cosim stage} --- implement the cosim backend the budget already reserves for.
  \item \textbf{Streaming/dataflow kernels} --- add an \texttt{hls::stream} task to
    exercise deadlock/invalid-streaming repair (\texttt{TIMEOUT\_OR\_DEADLOCK}).
  \item \textbf{Numerical-tolerance compare} --- honor a spec-declared tolerance in csim
    checking.
  \item \textbf{Compile-failure fixture} --- a fixture that proves the
    \texttt{compile\_error} $\rightarrow$ \texttt{minimal\_compile\_fix} path end-to-end.
\end{enumerate}

None of the items in this list change the paper's claims, which are scoped to what is
proven. The fidelity limitation above (Sec.~\ref{app:coverage:implverify}) is the one
that does bear on them, and is stated there rather than deferred; this section exists
to state both honestly.
